% GENERATED FILE. Edit canonical lessons, then run npm run generate:books.
% canonical-source-hash: fnv1a64:a4b982fa722c7346
% canonical-lessons: FR-C93-mince, FR-C93-gros, FR-C93-lourd, FR-C93-leger, FR-C93-sucre

\chapter{Thin, Big, Heavy, Light, Sweet}
\label{ch:fr-thin-big-heavy-light-sweet}
\hlchaptermodality{93}

\begin{chapteropening}
\textbf{By the end of this chapter:} \emph{I can name thin, big, heavy, light (not heavy), sweet (with sugar).}

Complete the last lesson of chapter 93: I can name thin, big, heavy, light (not heavy), sweet (with sugar).
\end{chapteropening}

\section[mince]{mince — thin, slim}

\label{lesson:FR-C93-mince}



\begin{quote}
Before the new one: say the French for kind, then the French for narrow.
\end{quote}

\subsection*{You'll want to know: mince}
\textbf{mince} — \textquotedblleft{}thin, slim\textquotedblright{}.

\begin{cousinweb}
\textbf{mince} — formed from the Old French verb \emph{mincier}, \textquotedblleft{}to cut fine\textquotedblright{}, from Vulgar Latin \emph{*minutiare}, built on Latin \emph{minutus}, \textquotedblleft{}small\textquotedblright{} (EtymDB 2.1 records the Latin word).
\end{cousinweb}

\begin{grammarlens}[title={What you've built}]
One more word for this chapter.
\end{grammarlens}

\subsection*{Your turn}
\begin{itemize}
\raggedright
  \item \emph{Say it:} \emph{mince}
  \item \emph{Say it:} \emph{mince}, once more
  \item \emph{Say it:} say \emph{mince}
  \item \emph{Recall:} say the French for kind, then the French for narrow, then say \emph{mince} again
\end{itemize}

\begin{tcolorbox}[breakable,colback=teal!4,colframe=teal!35!black,title={Before you move on}]
What does mince mean? (\textquotedblleft{}Thin, slim\textquotedblright{}.) Say it once more.
\end{tcolorbox}
\section[gros]{gros — big, fat}

\label{lesson:FR-C93-gros}



\begin{quote}
Before the new one: say the French for narrow, then the French for thin.
\end{quote}

\subsection*{You'll want to know: gros}
\textbf{gros} — \textquotedblleft{}big, fat\textquotedblright{}.

\begin{cousinweb}
\textbf{gros} $\leftarrow$ Late Latin \emph{grossus}, \textquotedblleft{}thick, coarse\textquotedblright{} — inherited. (EtymDB 2.1)
\end{cousinweb}

\begin{grammarlens}[title={What you've built}]
One more word for this chapter.
\end{grammarlens}

\subsection*{Your turn}
\begin{itemize}
\raggedright
  \item \emph{Say it:} \emph{gros}
  \item \emph{Say it:} \emph{gros}, once more
  \item \emph{Say it:} say \emph{gros}
  \item \emph{Recall:} say the French for narrow, then the French for thin, then say \emph{gros} again
\end{itemize}

\begin{tcolorbox}[breakable,colback=teal!4,colframe=teal!35!black,title={Before you move on}]
What does gros mean? (\textquotedblleft{}Big, fat\textquotedblright{}.) Say it once more.
\end{tcolorbox}
\section[lourd]{lourd — heavy}

\label{lesson:FR-C93-lourd}



\begin{quote}
Before the new one: say the French for thin, then the French for big.
\end{quote}

\subsection*{You'll want to know: lourd}
\textbf{lourd} — \textquotedblleft{}heavy\textquotedblright{}.

\begin{cousinweb}
\textbf{lourd} $\leftarrow$ Old French \emph{lourt}, \textquotedblleft{}dull, stupid\textquotedblright{}, usually traced to Latin \emph{luridus}, \textquotedblleft{}sallow, ghastly\textquotedblright{} — inherited. (EtymDB 2.1) The sense moved from dull to heavy.
\end{cousinweb}

\begin{grammarlens}[title={What you've built}]
One more word for this chapter.
\end{grammarlens}

\subsection*{Your turn}
\begin{itemize}
\raggedright
  \item \emph{Say it:} \emph{lourd}
  \item \emph{Say it:} \emph{lourd}, once more
  \item \emph{Say it:} say \emph{lourd}
  \item \emph{Recall:} say the French for thin, then the French for big, then say \emph{lourd} again
\end{itemize}

\begin{tcolorbox}[breakable,colback=teal!4,colframe=teal!35!black,title={Before you move on}]
What does lourd mean? (\textquotedblleft{}Heavy\textquotedblright{}.) Say it once more.
\end{tcolorbox}
\section[léger]{léger — light (not heavy)}

\label{lesson:FR-C93-leger}



\begin{quote}
Before the new one: say the French for big, then the French for heavy.
\end{quote}

\subsection*{You'll want to know: léger}
\textbf{léger} — \textquotedblleft{}light (not heavy)\textquotedblright{}.

\begin{cousinweb}
\textbf{léger} $\leftarrow$ Vulgar Latin \emph{*leviarius}, from Latin \emph{levis}, \textquotedblleft{}light\textquotedblright{} — inherited. (EtymDB 2.1)
\end{cousinweb}

\begin{grammarlens}[title={What you've built}]
One more word for this chapter.
\end{grammarlens}

\subsection*{Your turn}
\begin{itemize}
\raggedright
  \item \emph{Say it:} \emph{léger}
  \item \emph{Say it:} \emph{léger}, once more
  \item \emph{Say it:} say \emph{léger}
  \item \emph{Recall:} say the French for big, then the French for heavy, then say \emph{léger} again
\end{itemize}

\begin{tcolorbox}[breakable,colback=teal!4,colframe=teal!35!black,title={Before you move on}]
What does léger mean? (\textquotedblleft{}Light (not heavy)\textquotedblright{}.) Say it once more.
\end{tcolorbox}
\section[sucré]{sucré — sweet (with sugar)}

\label{lesson:FR-C93-sucre}



\begin{quote}
Before the new one: say the French for heavy, then the French for light.
\end{quote}

\subsection*{You'll want to know: sucré}
\textbf{sucré} — \textquotedblleft{}sweet (with sugar)\textquotedblright{}.

\begin{cousinweb}
\textbf{sucré} is the past participle of the verb for \textquotedblleft{}to sugar\textquotedblright{}, made from \emph{sucre}, \textquotedblleft{}sugar\textquotedblright{}: sugared.
\end{cousinweb}

\begin{grammarlens}[title={What you've built}]
One more word for this chapter.
\end{grammarlens}

\subsection*{Your turn}
\begin{itemize}
\raggedright
  \item \emph{Say it:} \emph{sucré}
  \item \emph{Say it:} \emph{sucré}, once more
  \item \emph{Say it:} say \emph{sucré}
  \item \emph{Recall:} say the French for heavy, then the French for light, then say \emph{sucré} again
\end{itemize}

\begin{tcolorbox}[breakable,colback=teal!4,colframe=teal!35!black,title={Before you move on}]
What does sucré mean? (\textquotedblleft{}Sweet (with sugar)\textquotedblright{}.) Say it once more.
\end{tcolorbox}
